% !TeX root = ../main.tex
\section{Краткое функциональное и математическое описание}

\subsection{Квадратурная демодуляция и фильтрация нижних частот}

Входной сигнал представляет собой сумму полезного ЛЧМ-радиоимпульса $s(t)$ и аддитивного белого гауссовского шума $n(t)$:
\begin{equation}
    x(t) = \operatorname{Re}\left\{ s_{\text{bb}}(t) \cdot e^{j 2\pi f_c t} \right\} + n(t),
\end{equation}
где $s_{\text{bb}}(t) = \exp\left(j \pi \frac{\Delta f}{T_{\text{и}}} t^2\right)$ при $|t| \le T_{\text{и}}/2$ --- комплексная огибающая симметричного ЛЧМ-импульса.

Аналого-цифровой преобразователь квантует сигнал с шагом $\Delta U = \frac{2 U_{\max}}{2^{n_a} - 1}$, формируя последовательность $x_{\text{adc}}[n]$. Цифровой гетеродин (NCO) генерирует отсчеты комплексной экспоненты:
\begin{equation}
    s_{\text{nco}}[n] = \exp\left(-j 2\pi \frac{f_G}{f_s} n\right).
\end{equation}

В результате перемножения формируется демодулированный сигнал:
\begin{equation}
    x_{\text{d}}[n] = x_{\text{adc}}[n] \cdot s_{\text{nco}}[n].
\end{equation}

Спектр полезного сигнала переносится в окрестность нулевой частоты ($f_c - f_G = 0$ МГц), а паразитная высокочастотная составляющая на суммарной частоте $f_c + f_G$ подавляется цифровым фильтром нижних частот (FIR LPF):
\begin{equation}
    x_{\text{fir}}[n] = \sum_{k=0}^{M} h_{\text{lpf}}[k] \cdot x_{\text{d}}[n-k],
\end{equation}
где $h_{\text{lpf}}[k]$ --- импульсная характеристика КИХ-фильтра порядка $M=128$.

\subsection{Математическая модель согласованного фильтра}

Импульсная характеристика согласованного фильтра $h[n]$ длительностью $N_{\text{имп}} = 4608$ отсчетов является комплексно-сопряженной и зеркально обращенной копией зондирующего видеосигнала:
\begin{equation}
    h[n] = s_{\text{bb}}^*[N_{\text{имп}} - 1 - n], \quad n = 0, 1, \dots, N_{\text{имп}} - 1.
\end{equation}

Выходной сигнал согласованного фильтра формируется линейной сверткой:
\begin{equation}
    x_{\text{mf}}[n] = \sum_{k=0}^{N_{\text{имп}}-1} x_{\text{fir}}[n - k] \cdot h[k].
\end{equation}

В работе реализовано три эквивалентных метода вычисления:
\begin{enumerate}
    \item \textbf{Прямая свертка (FIR):} пошаговое вычисление свертки во временной области с вычислительной сложностью $\mathcal{O}(L \cdot N_{\text{имп}})$.
    \item \textbf{Спектральный метод (FFT/IFFT):} быстрое вычисление циклической свертки через БПФ с дополнением нулями до размера $N_{\text{fft}} \ge L + N_{\text{имп}} - 1$:
    \begin{equation}
        x_{\text{mf}} = \operatorname{IFFT}\Big( \operatorname{FFT}(x_{\text{fir}}) \odot \operatorname{FFT}(h) \Big).
    \end{equation}
    \item \textbf{Сегментированная быстрая свертка (fftfilt):} потоковое вычисление линейной свертки методом перекрытия с накоплением (Overlap-Add).
\end{enumerate}

Для снижения уровня боковых лепестков импульсная характеристика модифицируется с помощью весового окна Хэмминга:
\begin{equation}
    h_w[n] = h[n] \cdot \left[ 0.54 - 0.46 \cos\left(\frac{2\pi n}{N_{\text{имп}} - 1}\right) \right].
\end{equation}
